% Options for packages loaded elsewhere
\PassOptionsToPackage{unicode}{hyperref}
\PassOptionsToPackage{hyphens}{url}
\documentclass[
]{article}
\usepackage{xcolor}
\usepackage[margin=1in]{geometry}
\usepackage{amsmath,amssymb}
\setcounter{secnumdepth}{-\maxdimen} % remove section numbering
\usepackage{iftex}
\ifPDFTeX
  \usepackage[T1]{fontenc}
  \usepackage[utf8]{inputenc}
  \usepackage{textcomp} % provide euro and other symbols
\else % if luatex or xetex
  \usepackage{unicode-math} % this also loads fontspec
  \defaultfontfeatures{Scale=MatchLowercase}
  \defaultfontfeatures[\rmfamily]{Ligatures=TeX,Scale=1}
\fi
\usepackage{lmodern}
\ifPDFTeX\else
  % xetex/luatex font selection
\fi
% Use upquote if available, for straight quotes in verbatim environments
\IfFileExists{upquote.sty}{\usepackage{upquote}}{}
\IfFileExists{microtype.sty}{% use microtype if available
  \usepackage[]{microtype}
  \UseMicrotypeSet[protrusion]{basicmath} % disable protrusion for tt fonts
}{}
\makeatletter
\@ifundefined{KOMAClassName}{% if non-KOMA class
  \IfFileExists{parskip.sty}{%
    \usepackage{parskip}
  }{% else
    \setlength{\parindent}{0pt}
    \setlength{\parskip}{6pt plus 2pt minus 1pt}}
}{% if KOMA class
  \KOMAoptions{parskip=half}}
\makeatother
\usepackage{color}
\usepackage{fancyvrb}
\newcommand{\VerbBar}{|}
\newcommand{\VERB}{\Verb[commandchars=\\\{\}]}
\DefineVerbatimEnvironment{Highlighting}{Verbatim}{commandchars=\\\{\}}
% Add ',fontsize=\small' for more characters per line
\usepackage{framed}
\definecolor{shadecolor}{RGB}{248,248,248}
\newenvironment{Shaded}{\begin{snugshade}}{\end{snugshade}}
\newcommand{\AlertTok}[1]{\textcolor[rgb]{0.94,0.16,0.16}{#1}}
\newcommand{\AnnotationTok}[1]{\textcolor[rgb]{0.56,0.35,0.01}{\textbf{\textit{#1}}}}
\newcommand{\AttributeTok}[1]{\textcolor[rgb]{0.13,0.29,0.53}{#1}}
\newcommand{\BaseNTok}[1]{\textcolor[rgb]{0.00,0.00,0.81}{#1}}
\newcommand{\BuiltInTok}[1]{#1}
\newcommand{\CharTok}[1]{\textcolor[rgb]{0.31,0.60,0.02}{#1}}
\newcommand{\CommentTok}[1]{\textcolor[rgb]{0.56,0.35,0.01}{\textit{#1}}}
\newcommand{\CommentVarTok}[1]{\textcolor[rgb]{0.56,0.35,0.01}{\textbf{\textit{#1}}}}
\newcommand{\ConstantTok}[1]{\textcolor[rgb]{0.56,0.35,0.01}{#1}}
\newcommand{\ControlFlowTok}[1]{\textcolor[rgb]{0.13,0.29,0.53}{\textbf{#1}}}
\newcommand{\DataTypeTok}[1]{\textcolor[rgb]{0.13,0.29,0.53}{#1}}
\newcommand{\DecValTok}[1]{\textcolor[rgb]{0.00,0.00,0.81}{#1}}
\newcommand{\DocumentationTok}[1]{\textcolor[rgb]{0.56,0.35,0.01}{\textbf{\textit{#1}}}}
\newcommand{\ErrorTok}[1]{\textcolor[rgb]{0.64,0.00,0.00}{\textbf{#1}}}
\newcommand{\ExtensionTok}[1]{#1}
\newcommand{\FloatTok}[1]{\textcolor[rgb]{0.00,0.00,0.81}{#1}}
\newcommand{\FunctionTok}[1]{\textcolor[rgb]{0.13,0.29,0.53}{\textbf{#1}}}
\newcommand{\ImportTok}[1]{#1}
\newcommand{\InformationTok}[1]{\textcolor[rgb]{0.56,0.35,0.01}{\textbf{\textit{#1}}}}
\newcommand{\KeywordTok}[1]{\textcolor[rgb]{0.13,0.29,0.53}{\textbf{#1}}}
\newcommand{\NormalTok}[1]{#1}
\newcommand{\OperatorTok}[1]{\textcolor[rgb]{0.81,0.36,0.00}{\textbf{#1}}}
\newcommand{\OtherTok}[1]{\textcolor[rgb]{0.56,0.35,0.01}{#1}}
\newcommand{\PreprocessorTok}[1]{\textcolor[rgb]{0.56,0.35,0.01}{\textit{#1}}}
\newcommand{\RegionMarkerTok}[1]{#1}
\newcommand{\SpecialCharTok}[1]{\textcolor[rgb]{0.81,0.36,0.00}{\textbf{#1}}}
\newcommand{\SpecialStringTok}[1]{\textcolor[rgb]{0.31,0.60,0.02}{#1}}
\newcommand{\StringTok}[1]{\textcolor[rgb]{0.31,0.60,0.02}{#1}}
\newcommand{\VariableTok}[1]{\textcolor[rgb]{0.00,0.00,0.00}{#1}}
\newcommand{\VerbatimStringTok}[1]{\textcolor[rgb]{0.31,0.60,0.02}{#1}}
\newcommand{\WarningTok}[1]{\textcolor[rgb]{0.56,0.35,0.01}{\textbf{\textit{#1}}}}
\usepackage{graphicx}
\makeatletter
\newsavebox\pandoc@box
\newcommand*\pandocbounded[1]{% scales image to fit in text height/width
  \sbox\pandoc@box{#1}%
  \Gscale@div\@tempa{\textheight}{\dimexpr\ht\pandoc@box+\dp\pandoc@box\relax}%
  \Gscale@div\@tempb{\linewidth}{\wd\pandoc@box}%
  \ifdim\@tempb\p@<\@tempa\p@\let\@tempa\@tempb\fi% select the smaller of both
  \ifdim\@tempa\p@<\p@\scalebox{\@tempa}{\usebox\pandoc@box}%
  \else\usebox{\pandoc@box}%
  \fi%
}
% Set default figure placement to htbp
\def\fps@figure{htbp}
\makeatother
\setlength{\emergencystretch}{3em} % prevent overfull lines
\providecommand{\tightlist}{%
  \setlength{\itemsep}{0pt}\setlength{\parskip}{0pt}}
\usepackage{bookmark}
\IfFileExists{xurl.sty}{\usepackage{xurl}}{} % add URL line breaks if available
\urlstyle{same}
\hypersetup{
  pdftitle={Revisar-conjunto-de-datos.R},
  pdfauthor={TRABAJO},
  hidelinks,
  pdfcreator={LaTeX via pandoc}}

\title{Revisar-conjunto-de-datos.R}
\author{TRABAJO}
\date{2026-08-18}

\begin{document}
\maketitle

\begin{Shaded}
\begin{Highlighting}[]
\CommentTok{\# Blas Avitia Arellano}
\CommentTok{\# 2327168}
\CommentTok{\# 18/08/2026}

\CommentTok{\#Importar Datos{-}{-}{-}{-}}
\CommentTok{\# Función read.csv importa datos de excel a R}

\NormalTok{IE }\OtherTok{\textless{}{-}} \FunctionTok{read.csv}\NormalTok{(}\StringTok{"Vivero.csv"}\NormalTok{, }\AttributeTok{header =}\NormalTok{ T)}
\NormalTok{IE}\SpecialCharTok{$}\NormalTok{Tratamiento }\OtherTok{\textless{}{-}} \FunctionTok{as.factor}\NormalTok{(IE}\SpecialCharTok{$}\NormalTok{Tratamiento)}

\CommentTok{\# Revisar{-}{-}{-}{-}{-}}
\CommentTok{\# Revisar solo una porción de datos}
\FunctionTok{head}\NormalTok{(IE) }\CommentTok{\# Primeras filas de BD IE}
\end{Highlighting}
\end{Shaded}

\begin{verbatim}
##   planta   IE Tratamiento
## 1      1 0.80        Ctrl
## 2      2 0.66        Ctrl
## 3      3 0.65        Ctrl
## 4      4 0.87        Ctrl
## 5      5 0.63        Ctrl
## 6      6 0.94        Ctrl
\end{verbatim}

\begin{Shaded}
\begin{Highlighting}[]
\FunctionTok{tail}\NormalTok{(IE) }\CommentTok{\# Ultimas filas de BD IE}
\end{Highlighting}
\end{Shaded}

\begin{verbatim}
##    planta   IE Tratamiento
## 37     37 0.88        Fert
## 38     38 1.15        Fert
## 39     39 0.88        Fert
## 40     40 0.78        Fert
## 41     41 1.16        Fert
## 42     42 0.91        Fert
\end{verbatim}

\begin{Shaded}
\begin{Highlighting}[]
\CommentTok{\# Graficar{-}{-}{-}{-}}

\FunctionTok{hist}\NormalTok{(IE}\SpecialCharTok{$}\NormalTok{IE,}
     \AttributeTok{main =} \StringTok{""}\NormalTok{,}
     \AttributeTok{xlab =} \StringTok{"Índice"}\NormalTok{,}
     \AttributeTok{ylab =} \StringTok{"Frecuencia"}\NormalTok{,}
     \AttributeTok{ylim =} \FunctionTok{c}\NormalTok{(}\DecValTok{0}\NormalTok{,}\DecValTok{12}\NormalTok{),}
     \AttributeTok{col =} \StringTok{"lightblue"}\NormalTok{)}
\end{Highlighting}
\end{Shaded}

\pandocbounded{\includegraphics[keepaspectratio]{Revisar-conjunto-de-datos_files/Revisar-conjunto-de-datos_files/figure-latex/unnamed-chunk-1-1.pdf}}

\begin{Shaded}
\begin{Highlighting}[]
\CommentTok{\# Normalidad datos {-}{-}{-}{-}}

\FunctionTok{shapiro.test}\NormalTok{(IE}\SpecialCharTok{$}\NormalTok{IE)}
\end{Highlighting}
\end{Shaded}

\begin{verbatim}
## 
##  Shapiro-Wilk normality test
## 
## data:  IE$IE
## W = 0.96225, p-value = 0.1777
\end{verbatim}

\begin{Shaded}
\begin{Highlighting}[]
\FunctionTok{mean}\NormalTok{(IE}\SpecialCharTok{$}\NormalTok{IE)}
\end{Highlighting}
\end{Shaded}

\begin{verbatim}
## [1] 0.8371429
\end{verbatim}

\begin{Shaded}
\begin{Highlighting}[]
\FunctionTok{sd}\NormalTok{(IE}\SpecialCharTok{$}\NormalTok{IE)}
\end{Highlighting}
\end{Shaded}

\begin{verbatim}
## [1] 0.1650319
\end{verbatim}

\begin{Shaded}
\begin{Highlighting}[]
\FunctionTok{var}\NormalTok{(IE}\SpecialCharTok{$}\NormalTok{IE)}
\end{Highlighting}
\end{Shaded}

\begin{verbatim}
## [1] 0.02723554
\end{verbatim}

\begin{Shaded}
\begin{Highlighting}[]
\FunctionTok{fivenum}\NormalTok{(IE}\SpecialCharTok{$}\NormalTok{IE)}
\end{Highlighting}
\end{Shaded}

\begin{verbatim}
## [1] 0.550 0.700 0.795 0.940 1.160
\end{verbatim}

\begin{Shaded}
\begin{Highlighting}[]
\CommentTok{\# Boxplot {-}{-}{-}{-}}
\FunctionTok{boxplot}\NormalTok{(IE}\SpecialCharTok{$}\NormalTok{IE }\SpecialCharTok{\textasciitilde{}}\NormalTok{ IE}\SpecialCharTok{$}\NormalTok{Tratamiento,}
        \AttributeTok{col =} \StringTok{"red"}\NormalTok{,}
        \AttributeTok{xlab =} \StringTok{"Tratamientos"}\NormalTok{,}
        \AttributeTok{ylab =} \StringTok{"Índice"}\NormalTok{) }
\end{Highlighting}
\end{Shaded}

\pandocbounded{\includegraphics[keepaspectratio]{Revisar-conjunto-de-datos_files/Revisar-conjunto-de-datos_files/figure-latex/unnamed-chunk-1-2.pdf}}

\begin{Shaded}
\begin{Highlighting}[]
\CommentTok{\# Revisar homogeneidad de varianzas{-}{-}{-}{-}}
\FunctionTok{bartlett.test}\NormalTok{(IE}\SpecialCharTok{$}\NormalTok{IE }\SpecialCharTok{\textasciitilde{}}\NormalTok{ IE}\SpecialCharTok{$}\NormalTok{Tratamiento)}
\end{Highlighting}
\end{Shaded}

\begin{verbatim}
## 
##  Bartlett test of homogeneity of variances
## 
## data:  IE$IE by IE$Tratamiento
## Bartlett's K-squared = 3.7423, df = 1, p-value = 0.05305
\end{verbatim}

\end{document}
